\begin{table}[htbp]
    \caption{Limites do valor ótimo de $Z$ por conjunto --- instância média}
    \label{tab:res-limites-medio}
    \centering
    \begin{tabular}{crrrr}
        \hline
        & \multicolumn{3}{c}{\textbf{Gurobi}} & \textbf{Limite combinatório} \\
        \cline{2-4} \cline{5-5}
        \textbf{Conjunto} & \textbf{Limite superior} & \textbf{Limite inferior} & \textbf{Gap} & \textbf{$LB_{comb}$} \\
        \hline
        01 & 79\,300,0 & 22\,350,0 & 71,8\% & 21\,800,0 \\
        02 & 73\,625,0 & 18\,000,0 & 75,6\% & 17\,500,0 \\
        03 & 91\,800,0 & 27\,400,0 & 70,2\% & 27\,200,0 \\
        04 & 84\,525,0 & 23\,800,0 & 71,8\% & 23\,500,0 \\
        05 & 89\,950,0 & 20\,450,0 & 77,3\% & 19\,900,0 \\
        06 & 70\,050,0 & 24\,550,0 & 65,0\% & 24\,400,0 \\
        07 & 83\,950,0 & 22\,750,0 & 72,9\% & 22\,400,0 \\
        08 & 88\,275,0 & 16\,500,0 & 81,3\% & 15\,900,0 \\
        09 & 92\,650,0 & 29\,400,0 & 68,3\% & 28\,700,0 \\
        10 & 86\,775,0 & 34\,300,0 & 60,5\% & 34\,200,0 \\
        \hline
    \end{tabular}
    \par\smallskip
    {\footnotesize Nota - Gap na convenção do \textit{solver}. $LB_{comb}$: limite inferior combinatório, conforme a Equação~\eqref{eq:limite-combinatorio}.}
\end{table}

\begin{table}[htbp]
    \caption{Resultados da heurística por conjunto --- instância média (10 execuções por conjunto; $Z$ sem a penalidade de capacidade)}
    \label{tab:res-z-medio}
    \centering
    \begin{adjustbox}{max width=\textwidth}
    \begin{tabular}{crrrrrcr}
        \hline
        \textbf{Conjunto} & \textbf{$Z$ da execução de menor $F$} & \textbf{$V_{cap}$} & \textbf{$Z$ médio} & \textbf{CV} & \textbf{Melhor $Z$ viável} & \textbf{Execuções viáveis} & \textbf{Excesso máx.} \\
        \hline
        01 & 28\,775,0 & 0 & 29\,777,5 & 2,5\% & 28\,775,0 & 10/10 & 28,7\% \\
        02 & 22\,750,0 & 0 & 23\,307,5 & 1,5\% & 22\,750,0 & 10/10 & 26,4\% \\
        03 & 32\,200,0 & 0 & 32\,865,0 & 1,7\% & 32\,200,0 & 10/10 & 17,5\% \\
        04 & 27\,825,0 & 0 & 28\,560,0 & 1,9\% & 27\,825,0 & 10/10 & 16,9\% \\
        05 & 25\,525,0 & 0 & 26\,387,5 & 2,5\% & 25\,525,0 & 10/10 & 24,8\% \\
        06 & 29\,550,0 & 0 & 30\,192,5 & 2,2\% & 29\,550,0 & 10/10 & 20,4\% \\
        07 & 27\,375,0 & 0 & 28\,425,0 & 2,4\% & 27\,375,0 & 10/10 & 20,3\% \\
        08 & 20\,050,0 & 0 & 21\,042,5 & 3,2\% & 20\,050,0 & 10/10 & 21,5\% \\
        09 & 33\,050,0 & 0 & 34\,067,5 & 1,4\% & 33\,050,0 & 10/10 & 12,4\% \\
        10 & 39\,825,0 & 0 & 40\,485,0 & 1,3\% & 39\,825,0 & 10/10 & 16,1\% \\
        \hline
    \end{tabular}
    \end{adjustbox}
    \par\smallskip
    {\footnotesize Nota - $Z$ da execução de menor $F$: valor de $Z$ da execução com menor \textit{fitness}, conforme a Equação~\eqref{eq:fitness}. $V_{cap}$: violação de capacidade dessa execução, em pacientes-dia, conforme a Equação~\eqref{eq:violacao-capacidade}, sendo $V_{cap} > 0$ uma solução inviável para o modelo exato. CV: coeficiente de variação de $Z$ entre as execuções (desvio padrão dividido pela média), como medida de dispersão. Melhor $Z$ viável: menor $Z$ entre as soluções viáveis encontradas nas execuções (--- se nenhuma). Excesso máx.: excesso máximo da execução de menor $F$ sobre o ótimo, conforme a Equação~\eqref{eq:excesso-otimo}, calculado com o maior entre o limite inferior do Gurobi e $LB_{comb}$.}
\end{table}
